\textbf{Решение задачи Д-1}

\vspace{1em}
\textbf{1. Исходные данные и перевод единиц}

Из условия задачи известны следующие параметры:
\begin{itemize}
    \item Шумовая температура антенны: $T_A = 100 \text{ К}$
    \item Коэффициент модуляции АМ сигнала: $M = 100\% = 1$
    \item Эффективное напряжение несущей: $U_{\text{нес.эфф}} = 1.6 \text{ мкВ} = 1.6 \cdot 10^{-6} \text{ В}$
    \item Входное сопротивление БВЧ: $R_{\text{вх}} = 75 \text{ Ом}$
    \item Коэффициент шума БВЧ: $F_{\text{дБ}} = 6 \text{ дБ}$
    \item Шумовая полоса пропускания: $\Pi_{\text{ш}} = 8 \text{ МГц} = 8 \cdot 10^6 \text{ Гц}$
\end{itemize}
Постоянная Больцмана: $k = 1.38 \cdot 10^{-23} \text{ Вт/(Гц}\cdot\text{К)}$. 
Стандартная температура (по умолчанию): $T_0 = 290 \text{ К}$.

\vspace{1em}
\textbf{2. Расчет мощности сигнала и шума на входе детектора}

Определим мощность несущей полезного сигнала на входе:
\begin{equation*}
    P_c = \frac{U_{\text{нес.эфф}}^2}{R_{\text{вх}}} = \frac{(1.6 \cdot 10^{-6})^2}{75} = \frac{2.56 \cdot 10^{-12}}{75} \approx 3.413 \cdot 10^{-14} \text{ Вт}
\end{equation*}

Переведем коэффициент шума БВЧ из децибел в разы:
\begin{equation*}
    F = 10^{\frac{6}{10}} = 10^{0.6} \approx 3.98 \text{ раза}
\end{equation*}

Рассчитаем эквивалентную шумовую температуру системы (приведенную ко входу):
\begin{equation*}
    T_{\Sigma} = T_A + T_0(F - 1) = 100 + 290(3.98 - 1) = 100 + 864.2 = 964.2 \text{ К}
\end{equation*}

Мощность шума в полосе БВЧ на входе детектора:
\begin{equation*}
    P_{\text{ш}} = k T_{\Sigma} \Pi_{\text{ш}} = 1.38 \cdot 10^{-23} \cdot 964.2 \cdot 8 \cdot 10^6 \approx 1.065 \cdot 10^{-13} \text{ Вт}
\end{equation*}

Вычислим отношение сигнал/шум (по мощности) на входе детектора:
\begin{equation*}
    q_{\text{вх}}^2 = \frac{P_c}{P_{\text{ш}}} = \frac{3.413 \cdot 10^{-14}}{1.065 \cdot 10^{-13}} \approx 0.32
\end{equation*}

\vspace{1em}
\textbf{3. Сравнение диодного амплитудного и синхронного детекторов}

Так как отношение сигнал/шум на входе $q_{\text{вх}}^2 \approx 0.32 < 1$, линейный диодный амплитудный детектор (детектор огибающей) будет работать в режиме \textbf{слабого сигнала}. В этом режиме проявляется ярко выраженный \textbf{пороговый эффект} — подавление полезного сигнала шумом. Отношение сигнал/шум на его выходе становится пропорциональным квадрату входного отношения:
\begin{equation*}
    q_{\text{вых, АД}}^2 \sim q_{\text{вх}}^4
\end{equation*}

Синхронный же детектор является линейным параметрическим устройством при любых отношениях сигнал/шум, и пороговый эффект в нем отсутствует. Отношение сигнал/шум на его выходе строго пропорционально входному:
\begin{equation*}
    q_{\text{вых, СД}}^2 \sim q_{\text{вх}}^2
\end{equation*}

\textbf{Вывод 1:} Да, при замене диодного детектора на синхронный отношение сигнал/шум на выходе \textbf{увеличится} (помехоустойчивость возрастет).

\vspace{1em}
\textbf{4. Оценка выигрыша в децибелах}

Отношение (выигрыш) мощностей сигнал/шум на выходе при использовании синхронного детектора вместо амплитудного оценивается коэффициентом подавления $K_{\text{под}} \approx q_{\text{вх}}^2$ (для слабых сигналов):
\begin{equation*}
    B = \frac{q_{\text{вых, СД}}^2}{q_{\text{вых, АД}}^2} \approx \frac{1}{q_{\text{вх}}^2}
\end{equation*}

Подставляем $q_{\text{вх}}^2 = 0.32$:
\begin{equation*}
    B = \frac{1}{0.32} = 3.125 \text{ раз}
\end{equation*}

Переведем полученный выигрыш в децибелы:
\begin{equation*}
    \Delta_{\text{дБ}} = 10 \log_{10}(B) = 10 \log_{10}(3.125) \approx 4.95 \text{ дБ}
\end{equation*}

\textit{Примечание: Если использовать более точную теоретическую формулу $B = \frac{1+2q_{\text{вх}}^2}{q_{\text{вх}}^2} = \frac{1}{q_{\text{вх}}^2} + 2$, то выигрыш составит $B = 3.125 + 2 = 5.125 \text{ раз}$, что соответствует $\approx 7.1 \text{ дБ}$. В инженерной практике чаще всего ограничиваются оценкой по главному подавляющему множителю (около 5 дБ).}

\textbf{Ответ:} Да, отношение сигнал/шум улучшится. Выигрыш составит приблизительно 5 дБ (по точной формуле $\approx 7.1$ дБ).